\section{局限性、治理策略与下一步研究议程}

Zou 在结尾强调了局限性，这部分非常关键，因为它决定系统能否进入真实科研生产。我们可以把局限分成能力边界、上下文边界、验证边界三类。

\subsection{能力边界：擅长工具使用，不擅长原创方法发明}

课程观点是当前 Agent 更擅长调用与组合既有工具，而非从零发明下一代基础方法。人类研究者在原创方法层仍有明显优势。两者互补路径是：人类造新方法，Agent 负责大规模扩散和落地应用。

\subsection{上下文边界：缺元数据时会做出无效假设}

Zou 给出一个典型问题：给 Agent 数据却不说明预处理流程，Agent 会自行补全前提并造成错误结论。因此，科学 Agent 要求更高的数据与流程元数据质量，不能仅靠``自然语言说明''。

\begin{warningbox}{科学场景下最危险的不是语法错误，而是前提错误}
前提错了，后续每一步都可能形式正确但语义失真。科研系统必须优先记录数据来源、预处理、实验条件、统计假设，并把这些信息纳入 Agent 的硬约束上下文。
\end{warningbox}

\subsection{验证边界：计算结论必须闭环到真实实验}

课程反复强调，很多科学问题最终要靠实验和现实世界数据验证。纯计算闭环最多给出候选，不应直接替代实验决策。稳健流程应是``Agent 生成候选 -> 人类筛选 -> 实验验证 -> 结果回流''。

\begin{knowledgebox}{面向实验科学的闭环模板}
\begin{itemize}
    \item 计算阶段：Agent 生成候选与置信区间。
    \item 评审阶段：人类确认可行性、伦理与资源约束。
    \item 实验阶段：按预注册协议执行验证。
    \item 回流阶段：把失败样本和异常元数据写回记忆库。
\end{itemize}
\end{knowledgebox}

\subsection{治理建议：从能力竞赛转向责任架构}

如果未来 AI-first 科研成为常态，治理重点应放在责任可追踪，而非只看结果新颖度。具体可从三层建设：过程可审计（日志、版本、调用链）、结果可核验（代码/数据/引用自动检查）、角色可问责（哪一步由谁主导、谁批准）。

\begin{importantbox}{可执行治理清单}
\begin{itemize}
    \item 对所有 Agent 调用保留结构化 trace 与配置快照。
    \item 将 reproducibility 检查与引用核验作为投稿前门槛。
    \item 对高风险结论设置强制 human sign-off。
    \item 在评审体系中纳入``协作质量''而非只看结论文本。
\end{itemize}
\end{importantbox}

\subsection{本章小结}

当前科学 Agent 的问题不是``能不能生成答案''，而是``答案如何被约束、验证并承担责任''。这决定了它能否从 demo 走向科学共同体的可信基础设施。

